\begin{frame}
  \frametitle{Z Transform and Convolution}
  
  Here's a formula for convolution:
  \[
  y[n] = \sum_{m=-\infty}^\infty h[m]x[n-m]
  \]
  Since the Z-transform of $x[n-m]$ is $z^{-m}X(z)$, we can transform
  the above difference equation to get
  \begin{align*}
  Y(z) &= \sum_{m=-\infty}^\infty h[m]z^{-m}X(z)\\
  &= H(z)X(z)
  \end{align*}
  So we confirm that $x[n]\ast h[n]\leftrightarrow H(z)X(z)$.
\end{frame}



%%%%%%%%%%%%%%%%%%%%%%%%%%%%%%%%%%%%%%%%%%%%
\section[Impulses]{Every Signal is a Weighted Sum of Impulses}
\setcounter{subsection}{1}

\begin{frame}
  \frametitle{Definition of the DTFT}

  The definition of the DTFT is
  \[
  X(\omega)  = \sum_{n=-\infty}^\infty x[n]e^{-j\omega n}
  \]
  \begin{itemize}
  \item We viewed this, before, as computing the phasor $X(\omega)$ using
    the orthogonality principle: multiply $x[n]$ by a pure tone at the
    corresponding frequency, and sum over all time.
  \item But there's another useful way to think about this.
  \end{itemize}
\end{frame}

\begin{frame}
  \frametitle{A Signal is a Weighted Sum of Impulses}

  Try writing the DTFT as
  \[
  X(\omega) = \ldots + x[-1]e^{j\omega} + x[0] + x[1]e^{-j\omega} + \ldots
  \]
  This looks like the DTFT of a difference equation.  The inverse DTFT would be
  \[
  x[n] = \ldots + x[-1]\delta[n+1] + x[0]\delta[n] + x[1]\delta[n-1] + \ldots
  \]
\end{frame}
  
\begin{frame}
  \frametitle{A Signal is a Weighted Sum of Impulses}
  
  \centerline{\animategraphics[loop,controls,width=4.5in]{5}{exp/impulses}{0}{59}}

\end{frame}

\begin{frame}
  \frametitle{A Signal is a Weighted Sum of Impulses}

  So we can use the DTFT formula, $X(\omega)=\sum x[n]e^{-j\omega n}$,
  to inspire us to think about $x[n]$ as just a weighted sum of impulses:

  \vspace*{1cm}
  
  \[
  X(\omega)=\sum_{m=-\infty}^\infty x[m]e^{-j\omega m}
  \leftrightarrow x[n]=\sum_{m=-\infty}^\infty x[m]\delta[n-m]
  \]
\end{frame}

